\documentclass[a4paper,12pt]{article}

\usepackage[utf8]{inputenc}
\usepackage[T1]{fontenc}
\usepackage[brazil]{babel}
\usepackage{geometry}
\usepackage{hyperref}
\usepackage{listings}
\usepackage{array}

\geometry{a4paper, margin=2.5cm}

\title{Relatório -- Palíndromo em DNA}
\author{}
\date{}

\begin{document}

\maketitle

\section{Ambiente}

Os testes foram executados em um ambiente Unix no Windows (WSL2), onde tenho meu workspace configurado para programação em C/C++. Para a etapa de compilação, utilizo o compilador G++, com o seguinte comando:

\begin{lstlisting}[language=bash]
g++-15 -std=gnu++23 -Wall -Wextra main.cpp -o program
\end{lstlisting}

\section{Desafios Enfrentados}

Sou um pouco metódico na hora de escrever código. Não penso apenas em desempenho, mas também me preocupo com a compreensão do código. Por isso, sempre procuro usar nomes de variáveis que condizem com o contexto da aplicação. Como Biologia Computacional não é uma área que domino, foi um pouco difícil decidir alguns nomes de forma que fizessem sentido no contexto do problema.

Além disso, a implementação do algoritmo para encontrar os palíndromos foi desafiadora. A princípio, pensei em adaptar o algoritmo de máximo subarray para o problema, mas percebi que não era o melhor caminho. Após pesquisar um pouco mais, encontrei o algoritmo
\href{https://www.geeksforgeeks.org/dsa/longest-palindromic-substring/}{Expand Around Center}, que se mostrou uma boa abordagem e me permitiu concluir o exercício com sucesso.

\section{Respostas às Perguntas}

\subsection*{2) Qual o tamanho do maior palíndromo maximal encontrado?}

O maior palíndromo maximal encontrado possui 16 bases. Ele foi identificado quando $k = 16$, na sequência de teste, como:

\begin{center}
\texttt{TATAAAACGTTTTATA}
\end{center}

Esse trecho começa na posição 23 da segunda linha do arquivo \texttt{dataset.txt}. A mesma sequência aparece novamente na terceira linha, começando na posição 90.

\subsection*{3) Enzimas de restrição}

Foram identificadas seis enzimas de restrição diferentes nos palíndromos encontrados nas sequências pesquisadas:

\begin{table}[h]
\centering
\caption{Enzimas de restrição identificadas.}
\begin{tabular}{|c|c|c|}
\hline
\textbf{Enzima} & \textbf{Sítio de reconhecimento} & \textbf{Palíndromo} \\
\hline
MseI  & \texttt{TTAA} & \texttt{TTAA} \\
\hline
TaqI  & \texttt{TCGA} & \texttt{TCGA} \\
\hline
DpnII & \texttt{GATC} & \texttt{GATC} \\
\hline
BstUI & \texttt{CGCG} & \texttt{CGCG} \\
\hline
AluI  & \texttt{AGCT} & \texttt{AGCT} \\
\hline
HhaI  & \texttt{GCGC} & \texttt{GCGC} \\
\hline
\end{tabular}
\end{table}

Fonte dos nomes das enzimas e dos respectivos sítios de reconhecimento:

\href{https://www.neb.com/en/tools-and-resources/selection-charts/alphabetized-list-of-recognition-specificities}
{NEB -- Alphabetized List of Recognition Specificities}.

\end{document}